\kchapter{Autoscaling}{hpa}

\begin{objectives}
\begin{itemize}
  \item the three dimensions of scaling: more Pods, bigger Pods, more nodes;
  \item the Horizontal Pod Autoscaler and its algorithm;
  \item the Vertical Pod Autoscaler and in-place resize;
  \item Cluster Autoscaler and Karpenter;
  \item event-driven scaling with KEDA.
\end{itemize}
\end{objectives}

\section{Three ways to scale}

\begin{center}
\resizebox{0.95\linewidth}{!}{%
\begin{tikzpicture}[
  card/.style={draw=k8sblue, fill=k8slight, rounded corners=6pt, line width=1pt,
               minimum width=4.6cm, minimum height=3.0cm, align=center, text width=4.2cm,
               font=\footnotesize, text=k8snavy}]
  \node[card] at (0,0) {{\large\bfseries\color{k8sblue}\faClone\ \ Horizontal}\\[4pt]
        \textbf{HPA} adds or removes Pod replicas.\\[3pt]
        \kimg[0.5cm]{pod}\kimg[0.5cm]{pod}\kimg[0.5cm]{pod} $\rightarrow$
        \kimg[0.5cm]{pod}\kimg[0.5cm]{pod}\kimg[0.5cm]{pod}\kimg[0.5cm]{pod}\kimg[0.5cm]{pod}};
  \node[card] at (5.2,0) {{\large\bfseries\color{k8sblue}\faIcon{expand-arrows-alt}\ \ Vertical}\\[4pt]
        \textbf{VPA} changes CPU/memory requests of each Pod.\\[3pt]
        \kimg[0.45cm]{pod} $\rightarrow$ \kimg[0.9cm]{pod}};
  \node[card] at (10.4,0) {{\large\bfseries\color{k8sblue}\faServer\ \ Cluster}\\[4pt]
        \textbf{Cluster Autoscaler / Karpenter} add or remove nodes.\\[3pt]
        \kimg[0.5cm]{node}\kimg[0.5cm]{node} $\rightarrow$ \kimg[0.5cm]{node}\kimg[0.5cm]{node}\kimg[0.5cm]{node}};
\end{tikzpicture}}
\end{center}

They work together: the HPA asks for more Pods; if there is no room, Pods go
\kc{Pending}; the Cluster Autoscaler sees Pending Pods and adds a node.

\section{Prerequisite: metrics-server}

The HPA needs live CPU/memory figures from \kterm{metrics-server}.

\begin{bashcode}{terminal}
kubectl apply -f https://github.com/kubernetes-sigs/metrics-server/releases/latest/download/components.yaml

# kind/local clusters use self-signed kubelet certs; allow them (lab only!)
kubectl patch deployment metrics-server -n kube-system --type=json \
  -p='[{"op":"add","path":"/spec/template/spec/containers/0/args/-","value":"--kubelet-insecure-tls"}]'

kubectl top nodes
kubectl top pods -A
\end{bashcode}

\section{Horizontal Pod Autoscaler}

\begin{yamlcode}{api-hpa.yaml}
apiVersion: autoscaling/v2
kind: HorizontalPodAutoscaler
metadata:
  name: api
spec:
  scaleTargetRef:
    apiVersion: apps/v1
    kind: Deployment
    name: api
  minReplicas: 2
  maxReplicas: 10
  metrics:
    - type: Resource
      resource:
        name: cpu
        target:
          type: Utilization
          averageUtilization: 70        # % of the CPU *request*
  behavior:
    scaleUp:
      stabilizationWindowSeconds: 0
      policies:
        - type: Percent
          value: 100                    # at most double every 30s
          periodSeconds: 30
    scaleDown:
      stabilizationWindowSeconds: 300   # wait 5 min before shrinking
\end{yamlcode}

\subsection{The algorithm}

Every 15 seconds the HPA computes:

\[
  \text{desiredReplicas} = \left\lceil \text{currentReplicas} \times
  \frac{\text{currentMetricValue}}{\text{targetMetricValue}} \right\rceil
\]

With 4 replicas at 140\,\% average CPU and a 70\,\% target:
\(\lceil 4 \times 140/70 \rceil = 8\) replicas.

\begin{warningbox}
Utilisation is measured against the \textbf{request}. A Deployment without CPU
requests cannot be autoscaled on CPU. Also, do not set \kc{replicas} in your
Deployment YAML when an HPA manages it --- every \kc{apply} would reset the
count.
\end{warningbox}

\begin{tryit}
\begin{bashcode}{terminal}
kubectl create deployment php --image=registry.k8s.io/hpa-example --port=80
kubectl set resources deployment php --requests=cpu=200m
kubectl expose deployment php --port=80
kubectl autoscale deployment php --cpu-percent=50 --min=1 --max=10

# generate load in a second terminal
kubectl run load --rm -it --image=busybox:1.36 -- \
  sh -c "while true; do wget -q -O- http://php; done"

kubectl get hpa php -w      # watch replicas climb, then fall after you stop the load
\end{bashcode}
\end{tryit}

\subsection{Custom and external metrics}

Through the \kc{custom.metrics.k8s.io} and \kc{external.metrics.k8s.io} APIs
(provided by adapters like the Prometheus Adapter), the HPA can scale on
requests per second, queue length, or any business metric:

\begin{yamlcode}{hpa-custom-metric.yaml}
metrics:
  - type: Pods
    pods:
      metric:
        name: http_requests_per_second
      target:
        type: AverageValue
        averageValue: "50"
\end{yamlcode}

\section{Vertical Pod Autoscaler}

The \kterm{VPA} (an add-on) watches real usage and recommends --- or applies ---
better requests. Start with \kc{updateMode: "Off"} to only collect
recommendations.

\begin{yamlcode}{api-vpa.yaml}
apiVersion: autoscaling.k8s.io/v1
kind: VerticalPodAutoscaler
metadata:
  name: api
spec:
  targetRef:
    apiVersion: apps/v1
    kind: Deployment
    name: api
  updatePolicy:
    updateMode: "Off"          # "Off" = recommend only; "Recreate"/"Auto" apply changes
\end{yamlcode}

\begin{bashcode}{terminal}
kubectl describe vpa api       # see "Target", "Lower Bound", "Upper Bound"
\end{bashcode}

Do not let the HPA and VPA both act on CPU/memory for the same workload: they
will fight. Recent Kubernetes versions also support \emph{in-place Pod
resize}, changing requests without restarting the Pod.

\section{Scaling the cluster itself}

\begin{itemize}
  \item \textbf{Cluster Autoscaler} adds nodes to a node group when Pods are
        Pending for lack of resources and removes under-used nodes.
  \item \textbf{Karpenter} (AWS, Azure, and others) provisions right-sized
        nodes on demand directly, often faster and cheaper.
\end{itemize}

\section{Event-driven autoscaling with KEDA}

\textbf{KEDA} scales workloads on external events --- Kafka lag, RabbitMQ queue
depth, SQS messages, cron schedules, Prometheus queries --- and can scale all
the way \emph{to zero} when there is no work.

\begin{yamlcode}{keda-scaledobject.yaml}
apiVersion: keda.sh/v1alpha1
kind: ScaledObject
metadata:
  name: worker
spec:
  scaleTargetRef:
    name: worker               # a Deployment
  minReplicaCount: 0
  maxReplicaCount: 30
  triggers:
    - type: rabbitmq
      metadata:
        queueName: emails
        mode: QueueLength
        value: "20"            # one replica per 20 waiting messages
      authenticationRef:
        name: rabbitmq-auth
\end{yamlcode}

\begin{recap}
\begin{itemize}
  \item HPA = more Pods; VPA = bigger Pods; Cluster Autoscaler/Karpenter = more nodes.
  \item HPA needs metrics-server and CPU/memory \emph{requests}.
  \item Tune \kc{behavior} to scale up fast and down slowly.
  \item KEDA scales on events and down to zero.
\end{itemize}
\end{recap}
